% ─────────────────────────────────────────────────────────────────────────────
\chapter{Conclusions, Limitations, and Future Work}
\label{chap:conclusions}
% ─────────────────────────────────────────────────────────────────────────────


% ─────────────────────────────────────────────────────────────────────────────
\section{Summary of the Research}
\label{sec:conclusions-summary}
% ─────────────────────────────────────────────────────────────────────────────

This thesis investigated the operational robustness of \gls{ai}-based systems
used for automated image classification and triage in forensic contexts. The
study was motivated by the observation that nominal performance on clean and
well-structured data does not fully characterize the behavior of systems that
may be applied to heterogeneous, degraded, distributionally different, or
deliberately manipulated digital evidence.

Within a \gls{df} workflow, classification errors may affect not only
technical performance but also the prioritization of material, the allocation
of analyst attention, and the effectiveness of the subsequent human review
process. A missed \texttt{weapon} image may fail to receive appropriate
priority, whereas an incorrect positive assignment may increase the volume of
irrelevant material requiring examination.

The research therefore adopted an operational perspective. Adversarial
attacks, anti-forensic transformations, and \gls{ood} inputs were treated as
distinct experimental conditions for examining system behavior beyond the
clean baseline. The objective was not to construct a conventional
\gls{advml} benchmark or to identify a universally superior system. Instead,
the study assessed whether nominal performance remained stable when the input
conditions changed and examined the resulting consequences for forensic
triage.

The experimental design combined two complementary evaluation levels. The
first comprised three transparent proxy models: \gls{efficientnetb0}, \gls{resnet18}, and
a \gls{clip}-based model consisting of a frozen visual encoder and a supervised
binary classification head. These models provided access to predicted classes,
maximum predicted-class probabilities, confusion matrices,
clean-to-perturbed error transitions, and internal gradients.

The second evaluation level comprised
\gls{magnetaxiomtool}/\gls{magnetai},
\gls{griffeye}/\gls{t3kcore},
\gls{excirefoto2025}, and
\gls{cellebriteinseyets}. These selected black-box software systems were
evaluated exclusively through the outputs observable to the analyst and through
explicitly documented normalization rules. The proxy models were not intended
to reproduce, approximate, or explain the proprietary internal models used by
the selected software systems.

A central methodological feature was the use of the same frozen source
population at both evaluation levels. The dataset comprised 1\,500 source
images: 500 \texttt{weapon} images, 500 \texttt{non\_weapon} images, and a
separate 500-image \gls{ood} branch. Adversarial and anti-forensic artifacts
were generated exclusively from the 1\,000-image binary subset.

The resulting forensic evaluation bundle contained 11\,500 logical items:
1\,000 clean binary inputs, 500 clean \gls{ood} inputs, 5\,000 adversarial
artifacts, and 5\,000 anti-forensic artifacts. The 10\,000 perturbed artifacts
remain derived versions of the same 1\,000 binary source images and do not
represent independent underlying cases.

The selected black-box configurations produced one normalized decision for
every bundle item, whereas the proxy-model pipeline preserved the fold-aware
evaluation design described in Chapter~\ref{chap:implementation}. This common
experimental basis made it possible to compare observable operating profiles
without assuming architectural or semantic equivalence among the evaluated
systems.

The complete procedure was formalized through the \gls{fairlab} protocol.
\gls{fairlab} is not a new classifier or standalone software platform. It is a
task-specific evaluation protocol connecting dataset construction, human
semantic validation, deterministic partitioning, proxy-model evaluation,
perturbation generation, blind-bundle construction, black-box output
normalization, quantitative metrics, qualitative explainability analysis, and
traceability controls.

Its purpose is to make operational robustness experimentally observable while
supporting procedural reproducibility and artifact-level auditability. The
protocol preserves the relationships among source images, derived artifacts,
model predictions, proprietary software outputs, normalized decisions, metrics,
and reporting assets.

The principal conclusion is that robustness cannot be inferred from clean
accuracy alone. It must be evaluated as a contextual property of a specific
system version, configuration, task definition, input population, output
mapping, and threat model.

Automated image-analysis outputs may provide useful support for forensic
triage, but they must remain traceable, reviewable, and subordinate to expert
assessment
\cite{costantini2019digital,sanna2024preliminary,sanna2025pixelperfect}.

% ─────────────────────────────────────────────────────────────────────────────
\section{Answers to the Research Questions}
\label{sec:conclusions-research-questions}
% ─────────────────────────────────────────────────────────────────────────────

\paragraph{\textbf{RQ1. To what extent do the transparent proxy models maintain
their clean-baseline performance under the separate \gls{ood}, adversarial,
and anti-forensic conditions?}}

The proxy models maintain their clean-baseline performance only to a
condition-dependent extent. All three architectures achieve comparatively
strong clean performance, with accuracies of 0.978 for \gls{efficientnetb0}, 0.964
for \gls{resnet18}, and 0.927 for the \gls{clip}-based proxy. These nominal results do
not remain uniformly representative under the other evaluation branches.

The five anti-forensic transformations generally produce smaller aggregate
accuracy variations than the strongest adversarial attacks. The largest
anti-forensic accuracy reduction among the proxy models is 0.030 for \gls{resnet18}
under \gls{histogrammodification}. Even where aggregate accuracy remains
unchanged or increases, however, clean-correct samples may become incorrect.
Anti-forensic robustness must therefore be interpreted through both aggregate
metrics and directional error transitions.

The direct-target adversarial conditions produce the most substantial
degradation. \gls{efficientnetb0} accuracy decreases from 0.978 to 0.575 under
\gls{fgsm} and to 0.015 under \gls{sigmazero}. The corresponding results show
that strong clean performance does not imply robustness to attacks optimized
against the evaluated model.

The separate \gls{ood} branch exposes an additional behavior that cannot be
predicted from clean binary accuracy. The fold-aggregated weapon assignment
rates are 0.3696 for \gls{efficientnetb0}, 0.4376 for \gls{resnet18}, and 0.8356 for the
\gls{clip}-based proxy. These values describe forced closed-set assignments and
not supervised correctness.

The proxy models therefore preserve nominal behavior under some frozen
conditions but not across the complete evaluation. Clean performance,
anti-forensic sensitivity, direct adversarial vulnerability, and \gls{ood}
assignment behavior remain distinct dimensions
(Sections~\ref{sec:results-clean-baseline}--\ref{sec:results-adversarial}).

\paragraph{\textbf{RQ2. How do direct source-model attacks,
cross-architecture transfer, and the model-agnostic \gls{colorshift} condition
differ in their effects on the evaluated proxy models?}}

The direct source-model attacks produce substantially larger effects than the
observed cross-architecture transfer. The four model-dependent attacks were
generated against fold-specific \gls{efficientnetb0} checkpoints. Their evaluation
on \gls{efficientnetb0} therefore represents direct-target assessment, whereas
their evaluation on \gls{resnet18} and the \gls{clip}-based proxy measures empirical
transfer.

The clearest direct-target failure is produced by \gls{sigmazero}, which
reduces \gls{efficientnetb0} accuracy to 0.015 and renders 963 of the 978
clean-correct samples incorrect. \gls{fgsm} reduces accuracy to 0.575 and
introduces 406 errors among clean-correct samples, predominantly in the
\texttt{weapon}-to-\texttt{non\_weapon} direction.

The same frozen artifacts produce more limited and heterogeneous effects on
the non-target architectures. The largest transferred degradation is observed
under \gls{fgsm}: \gls{resnet18} accuracy decreases from 0.964 to 0.918, whereas the
\gls{clip}-based proxy decreases from 0.927 to 0.914. Transferred
\gls{sigmazero} artifacts do not reproduce the direct-target collapse and
instead coexist with small aggregate accuracy increases and limited numbers of
new individual errors.

These results demonstrate that transferability must be measured rather than
assumed. Limited transfer does not establish direct robustness of \gls{resnet18} or
the \gls{clip}-based proxy because no matched attacks were optimized
independently against those architectures.

\gls{colorshift} follows a separate experimental rationale. It is deterministic
and model-agnostic and does not use gradients, parameters, scores, or
checkpoints. Its effects are comparatively limited and are interpreted as
condition-specific induced errors rather than direct or transferred attack
success
(Section~\ref{sec:results-adversarial}).

\paragraph{\textbf{RQ3. How do the selected black-box software systems respond
to the same frozen experimental corpus, and how do observable output semantics
and configuration choices affect the resulting operational profiles?}}

The selected black-box software systems produce heterogeneous operating
profiles even though they process the same frozen forensic evaluation bundle.
Their outputs are not internally equivalent: the evaluated decisions derive
from tags, bookmarks, exported classification labels, or semantic-retrieval
membership. The resulting metrics are therefore conditional on the explicit
normalization rule adopted for each configuration.

Under the frozen protocol, \gls{griffeye}/\gls{t3kcore} achieves a
bundle-weighted accuracy of 0.972 and an \gls{fpr} of 0.007. This selective
profile is inseparable from the conservative decision to treat only the
\texttt{CORE/Violence/Firearm} bookmark as a primary positive result.

\gls{cellebriteinseyets} produces a comparatively balanced profile, with
bundle accuracy, weapon recall, and weapon precision all close to 0.958,
together with similar overall false-negative and false-positive rates.

\gls{magnetaxiomtool}/\gls{magnetai} achieves a bundle accuracy of 0.933 and
an \gls{fnr} of 0.099. Its condition-level results show increased
missed-target exposure under transferred \gls{fgsm} artifacts and
\gls{gaussianblur}.

The three \gls{excirefoto2025} configurations provide the clearest evidence of
configuration sensitivity. On clean inputs, moving from \texttt{D20} to
\texttt{D80} reduces the \gls{fnr} from 0.120 to 0.014, while increasing the
\gls{fpr} from 0.026 to 0.156. The corresponding \gls{ood} weapon assignment
rate increases from 0.238 to 0.522. Under transferred \gls{sigmazero}
artifacts, the same change reduces the \gls{fnr} from 0.180 to 0.024 but
increases the \gls{fpr} from 0.190 to 0.662.

The results therefore do not support a universal ranking. More restrictive
configurations may reduce false-positive review workload while increasing
missed detections. More permissive configurations may reduce missed-target risk
while routing substantially more \texttt{non\_weapon} and \gls{ood} material
to the firearm-oriented review category.

Observable output semantics, category selection, retrieval settings, software
versions, enabled processing functions, and normalization rules are therefore
part of the evaluated operating profile
(Section~\ref{sec:results-commercial}).

\paragraph{\textbf{RQ4. Which error transitions, robustness degradations,
\gls{ood} assignment behaviors, and unknown output states are most relevant
to forensic triage and continued human review?}}

The most relevant error transition is
\texttt{weapon}\(\rightarrow\)\texttt{non\_weapon}, because it may prevent
potentially relevant firearm-oriented content from receiving appropriate review
priority. The opposite transition,
\texttt{non\_weapon}\(\rightarrow\)\texttt{weapon}, remains operationally
important because it increases analyst workload and may dilute prioritization.

Aggregate degradation must be interpreted jointly with these directions. The
\gls{efficientnetb0} \gls{sigmazero} result demonstrates a near-total direct-target
collapse involving both missed targets and increased false-positive workload.
Transferred \gls{fgsm} artifacts instead produce particularly relevant
false-negative burdens for Magnet.AI and Cellebrite Inseyets, whereas Excire
\texttt{D50} under transferred \gls{sigmazero} primarily exhibits an expansion
of false-positive semantic retrievals.

The \gls{ood} results identify a separate operational risk. An \gls{ood} input
assigned to \texttt{weapon} is not a supervised false positive, but it may
still be prioritized for review and increase workload. The evaluated proxy
models and all six black-box configurations produce non-zero weapon assignment
rates on the separate \gls{ood} branch.

Maximum predicted-class probabilities also require cautious interpretation.
The representative \gls{sigmazero} failure shows that an incorrect prediction
may be associated with a maximum probability of 1.000. Such a value does not
establish calibrated certainty, correctness, or forensic reliability.

No \texttt{unknown} normalized output occurs in the canonical 69\,000-decision
black-box result population. Consequently, unknown outputs were not an observed
failure mode in the frozen evaluation. Their explicit handling remains
methodologically necessary for future software versions, alternative exports,
or incomplete normalization cases.

The principal operational requirements are therefore systematic human review,
task-specific validation, documentation of software and model versions,
explicit recording of configurations and output mappings, preservation of the
links between source items and normalized decisions, and separate reporting of
missed-target risk, false-positive workload, and \gls{ood} assignment behavior
(Sections~\ref{sec:results-comparative} and
\ref{sec:results-limitations}).


% ─────────────────────────────────────────────────────────────────────────────
\section{Main Contributions}
\label{sec:conclusions-contributions}
% ─────────────────────────────────────────────────────────────────────────────

The contribution of this thesis is articulated across four closely connected
dimensions.

First, the thesis formalizes \gls{fairlab} as an integrated methodological
protocol for evaluating the operational robustness of \gls{ai}-based
image-analysis systems in \gls{df}. The protocol extends evaluation beyond
clean classification accuracy and connects human-validated dataset
construction, deterministic partitioning, transparent proxy-model analysis,
adversarial attacks, anti-forensic transformations, a separate \gls{ood}
branch, black-box software evaluation, output normalization, explainability,
and traceability controls.

The methodological contribution lies in the coordinated application of these
components within a single reproducible and auditable process. \gls{fairlab}
is not proposed as a new classifier, defensive mechanism, or standalone
forensic software platform. It provides a structured procedure through which
system behavior can be examined under controlled but operationally relevant
conditions.

Second, the work establishes a common frozen experimental basis for the
parallel evaluation of transparent proxy models and selected black-box software
systems. The source population contains three equally sized branches,
comprising 500 \texttt{weapon} images, 500 \texttt{non\_weapon} images, and
500 images reserved for the separate \gls{ood} analysis.

The 1\,000-image binary subset was used to generate five adversarial or
adversarial-style artifact populations and five anti-forensic artifact
populations. Together with the clean binary and \gls{ood} images, these
derivatives form the 11\,500-item forensic evaluation bundle.

The 10\,000 perturbed artifacts remain linked to their 1\,000 binary source
images and are not treated as independent underlying cases. This structure
allows clean and manipulated behavior to be compared while preserving the
identity and provenance of each source--condition combination.

Third, the thesis integrates transparent model evaluation with operational
black-box assessment without assuming architectural equivalence between the two
levels.

\gls{efficientnetb0}, \gls{resnet18}, and the \gls{clip}-based proxy provide inspectable
experimental references. They allow direct access to fold-specific
checkpoints, predicted classes, maximum predicted-class probabilities,
confusion matrices, matched clean-to-perturbed transitions, and internal
gradients.

The selected black-box systems instead expose heterogeneous operational
outputs. For \gls{magnetaxiomtool}/\gls{magnetai}, the observable output is the
\texttt{Possible weapons} tag, whereas
\gls{griffeye}/\gls{t3kcore} exposes the
\texttt{CORE/Violence/Firearm} bookmark.

The operational output adopted for \gls{excirefoto2025} is membership in the
union of eight firearm-oriented semantic-query result sets. For
\gls{cellebriteinseyets}, the evaluation uses the exported \texttt{Armi},
\texttt{Pistola}, and \texttt{Fucile} classification labels.

These outputs were normalized into a common operational
\texttt{weapon}/\texttt{non\_weapon} representation while keeping the
\gls{ood} branch analytically separate. The explicit normalization rules make
the comparison reproducible without implying that the selected systems
implement the same internal taxonomy or classification task.

The same model-dependent artifacts generated against fold-specific
\gls{efficientnetb0} checkpoints were subsequently evaluated on \gls{resnet18}, the
\gls{clip}-based proxy, and the selected black-box configurations. The protocol
therefore distinguishes direct-target vulnerability from empirical
cross-architecture and black-box transfer behavior. The model-agnostic
\gls{colorshift} condition and the anti-forensic transformations remain
separate from model-dependent attack success.

Fourth, the thesis provides an auditable and error-oriented research artifact.
Stable identifiers, source records, fold assignments, model checkpoints,
perturbation manifests, cryptographic digests, blind-bundle identifiers,
normalized predictions, metric tables, and reporting assets preserve the
relationship between each source image and the corresponding experimental
outcomes.

This traceability makes it possible to distinguish model or software behavior
from dataset-construction errors, missing outputs, normalization choices, and
residual experimental limitations. It also supports controlled reproducibility
where complete source data or proprietary software exports cannot be publicly
redistributed.

The operational analysis does not reduce robustness to a single aggregate
score. It examines false negatives affecting firearm-oriented content, false
positives increasing human-review workload, and directional
clean-to-perturbed error transitions.

The analysis also considers closed-set weapon assignments on the separate
\gls{ood} branch, configuration-dependent changes in recall and selectivity,
and the distinction between direct-target and transfer behavior under
adversarial conditions. Residual metadata sensitivity and the continued need
for human verification are incorporated into the resulting operational
interpretation.

A complementary qualitative \gls{xai} analysis was conducted using
\gls{ig} on five purposively selected \gls{efficientnetb0} cases. The resulting
visualizations support case-level diagnostic inspection but are not treated as
quantitative robustness evidence, semantic localization ground truth, or an
explanation of the selected black-box systems.

Overall, the contribution of the thesis lies not in proposing a universally
robust image classifier, but in defining and applying a traceable procedure for
revealing how heterogeneous automated systems behave when the assumptions of
clean, closed-set evaluation no longer hold.

% ─────────────────────────────────────────────────────────────────────────────
\section{Main Experimental Findings}
\label{sec:conclusions-findings}
% ─────────────────────────────────────────────────────────────────────────────

The results reported in Chapter~\ref{chap:results} show that nominal
performance and operational robustness are related but non-equivalent
properties. High clean accuracy provides an important baseline, but it does not
predict how a system will behave under distributional mismatch, image-level
transformations, adversarial manipulation, or alternative operational
configurations.

Table~\ref{tab:conclusions-main-findings} summarizes the principal findings and
their forensic relevance.

\begingroup
\small
\renewcommand{\arraystretch}{1.14}
\setlength{\tabcolsep}{4pt}
\begin{longtable}{@{}L{0.15\textwidth} L{0.39\textwidth} L{0.39\textwidth}@{}}
\caption{Summary of the main experimental findings}
\label{tab:conclusions-main-findings} \\
\toprule
\textbf{Area} &
\textbf{Main finding} &
\textbf{Operational relevance} \\
\midrule
\endfirsthead

\toprule
\textbf{Area} &
\textbf{Main finding} &
\textbf{Operational relevance} \\
\midrule
\endhead

\midrule
\multicolumn{3}{r}{\footnotesize Continued on the next page} \\
\endfoot

\bottomrule
\endlastfoot

Clean baseline
&
\gls{efficientnetb0} achieves the highest clean accuracy among the proxy models, at
0.978. Under the adopted conservative mapping,
\gls{griffeye}/\gls{t3kcore} achieves the highest clean accuracy among the
selected black-box configurations, also at 0.978.
\gls{cellebriteinseyets} produces slightly higher clean
\texttt{weapon} recall than \gls{griffeye}/\gls{t3kcore}.
&
Clean performance establishes the nominal reference but does not determine
behavior under \gls{ood}, adversarial, anti-forensic, or
configuration-dependent conditions. \\

\gls{ood}
&
The \gls{clip}-based proxy assigns 0.8356 of its fold-specific \gls{ood}
responses to \texttt{weapon}. \gls{efficientnetb0} and \gls{resnet18} produce lower weapon
assignment rates but substantially higher proportions of maximum predicted-class
probabilities above 0.90. The black-box weapon assignment rates range from
0.238 to 0.522 across the six configurations.
&
Inputs outside the supervised binary task may still be routed to the
firearm-oriented review category. These assignments are closed-set operating
behaviors rather than ordinary binary errors. \\

Anti-forensic
&
The five evaluated anti-forensic transformations generally produce smaller
aggregate accuracy variations than the strongest direct-target adversarial
conditions. Nevertheless, they introduce condition- and system-specific false
negatives and false positives. Gaussian Blur is particularly relevant for
\gls{efficientnetb0} and Magnet.AI, whereas Histogram Modification affects \gls{resnet18},
Griffeye/T3K CORE, Excire \texttt{D50}, and Cellebrite through different error
profiles.
&
Operationally plausible compression, resizing, blurring, and tonal
transformations may alter triage decisions even when the aggregate accuracy
change appears limited. Error direction remains more informative than aggregate
degradation alone. \\

Adversarial
&
\gls{efficientnetb0} undergoes severe direct-target degradation under
\gls{sigmazero}, with accuracy decreasing to 0.015, and under \gls{fgsm}, with
accuracy decreasing to 0.575. The same artifacts produce more limited and
heterogeneous transfer effects on \gls{resnet18}, the \gls{clip}-based proxy, and the
selected black-box configurations.
&
High clean accuracy does not guarantee direct adversarial robustness. Transfer
must be measured rather than assumed and must remain distinct from attacks
optimized directly against the evaluated system. \\

Black-box software
&
The selected configurations expose different operating profiles.
\gls{griffeye}/\gls{t3kcore} is highly selective under the conservative primary
bookmark mapping; \gls{cellebriteinseyets} presents a comparatively balanced
profile; \gls{magnetaxiomtool}/\gls{magnetai} shows greater missed-target
exposure under some conditions; and \gls{excirefoto2025} varies materially
across \texttt{D20}, \texttt{D50}, and \texttt{D80}.
&
No configuration minimizes false negatives, false positives, and \gls{ood}
weapon assignments simultaneously. Software version, enabled functions,
observable output semantics, configuration, and normalization mapping are part
of the reported result. \\

Metadata sensitivity
&
Fifteen of the 11\,500 blind-bundle files contain at least one predefined
semantically or experimentally sensitive metadata-term match. Excluding the
affected clean and \gls{ood} files produces only limited aggregate metric
variation.
&
The leave-out results support the stability of the aggregate conclusions but do
not establish that embedded metadata was unused or causally irrelevant to every
affected decision. \\

\gls{xai}
&
The qualitative \gls{ig} analysis examines one correct clean case, three
supervised failure cases, and one \gls{ood} closed-set assignment. The cases
include incorrect outputs associated with maximum predicted-class probabilities
of up to 1.000.
&
Explainability supports qualitative diagnosis of selected proxy-model outputs
but does not establish prediction correctness, calibrated certainty, causal
feature relevance, or black-box system behavior. \\

\end{longtable}
\endgroup

The clean results already reveal different operational error profiles.
\gls{efficientnetb0} minimizes the total number of clean errors among the proxy
models. The \gls{clip}-based proxy instead achieves the lowest clean
\gls{fnr}, at 0.012, while producing the highest clean \gls{fpr}, at 0.134.
A reduction in missed target images may therefore be accompanied by a
substantial increase in the amount of material requiring human review.

The separate \gls{ood} analysis identifies a dimension that cannot be inferred
from the clean binary results. The \gls{ood} population is heterogeneous and
includes borderline, degraded, synthetic, semantically adjacent, and otherwise
out-of-task content. It is not a third supervised class.

The observed weapon assignment rates show that both proxy models and selected
black-box configurations may route such content to the firearm-oriented
operational state. These results do not determine whether a specific
\gls{ood} image is relevant to an investigation. They characterize the
behavior of closed-set systems when their inputs do not conform to the nominal
binary task.

The maximum predicted-class probability results provide a separate diagnostic
dimension. The \gls{clip}-based proxy produces the highest \gls{ood} weapon
assignment rate while none of its fold-specific responses exceeds the 0.90
maximum-probability threshold. \gls{resnet18} produces fewer weapon assignments but
the highest high-maximum-probability rate.

The two indicators must therefore not be conflated. Assignment frequency does
not establish output concentration, and a high maximum predicted-class
probability does not establish correctness, calibration, or forensic
reliability.

The anti-forensic results show that model-agnostic transformations can affect
automated analysis without reproducing the extreme degradation observed under
the strongest direct-target adversarial attacks. Their operational importance
depends on the direction of the resulting errors.

For example, Gaussian Blur introduces 17 new \gls{efficientnetb0} errors, all in the
\texttt{weapon}-to-\texttt{non\_weapon} direction. Histogram Modification
introduces errors in both directions for \gls{resnet18}. At the black-box level,
Gaussian Blur produces the largest anti-forensic missed-target burden for
Magnet.AI, whereas Histogram Modification is associated with different
combinations of false negatives and false positives across Griffeye/T3K CORE,
Excire \texttt{D50}, and Cellebrite Inseyets.

A limited or negative accuracy drop therefore does not imply behavioral
invariance. Newly introduced errors may be offset by clean-condition errors
that become correct under the transformed population.

The adversarial analysis demonstrates the clearest separation between nominal
performance and direct-target robustness. The four model-dependent attacks were
generated against fold-specific \gls{efficientnetb0} checkpoints.
\gls{fgsm}, \gls{sigmazero}, and \gls{superdeepfool} used white-box access,
whereas \gls{onepixel} used score-based access.

Under direct-target evaluation, \gls{sigmazero} renders 963 of the 978
clean-correct \gls{efficientnetb0} inputs incorrect. The induced transitions are
nearly balanced between missed target images and false-positive assignments.
\gls{fgsm} produces fewer errors but a stronger concentration in the
\texttt{weapon}-to-\texttt{non\_weapon} direction.

The same artifacts do not reproduce equivalent effects on the non-target
architectures or selected black-box configurations. Their evaluation on those
systems measures empirical transfer rather than direct robustness. The more
limited effects must not be generalized as resistance to attacks optimized
independently against each evaluated architecture or commercial system.

The black-box findings reinforce the importance of configuration and output
semantics. The behavior of \gls{excirefoto2025} changes substantially between
\texttt{D20}, \texttt{D50}, and \texttt{D80}. Increasing the semantic-distance
limit reduces missed detections while increasing false-positive workload and
\gls{ood} weapon assignments.

The low \gls{fpr} observed for \gls{griffeye}/\gls{t3kcore} is conditional on
the conservative use of only \texttt{CORE/Violence/Firearm} as a primary
positive bookmark. The Cellebrite results depend on the selected exported
classification labels, whereas the Magnet.AI results depend on the observable
\texttt{Possible weapons} tag.

The resulting profiles are therefore properties of frozen system
configurations and normalization rules rather than intrinsic, version-independent
characteristics of the software families.

The embedded-metadata analysis identifies a residual limitation of blind-bundle
construction. Opaque file names, paths, and directories prevent direct
file-system disclosure of experimental assignments, but they do not remove
semantic information embedded within the image bytes.

The leave-out analysis shows that excluding the 11 affected clean images and
four affected \gls{ood} images does not materially alter the aggregate metrics.
It cannot determine whether any system accessed or relied on those metadata
fields. A causal assessment would require a paired metadata-neutralized bundle.

The five \gls{xai} cases provide qualitative diagnostic examples rather than a
representative statistical sample. The clean false negative and the
anti-forensic false negative illustrate missed-target behavior under nominal
and transformed conditions. The adversarial false positive demonstrates that an
incorrect output may be associated with a maximum predicted-class probability
of 1.000.

The \gls{ood} case instead documents a closed-set assignment outside the
supervised binary task and must not be described as an ordinary false positive.
The attribution maps represent attribution magnitudes only for the selected
\gls{efficientnetb0} output under the adopted attribution target, baseline, and
integration procedure.

Overall, the findings show that validation on clean data is not equivalent to
validation under operational conditions. Robustness assessment must jointly
consider baseline performance, error direction, \gls{ood} behavior,
adversarial and anti-forensic sensitivity, configuration dependence, output
mapping, traceability, and the continued role of human review.


% ─────────────────────────────────────────────────────────────────────────────
\section{Operational and Forensic Implications}
\label{sec:conclusions-operational-implications}
% ─────────────────────────────────────────────────────────────────────────────

The experimental evidence supports the use of automated image-analysis systems
as traceable triage aids rather than as autonomous mechanisms for evidentiary
decision-making. Their principal operational value lies in assisting analysts
with the prioritization and organization of large multimedia collections. This
value depends, however, on the possibility of reconstructing how an output was
produced, verifying it against the source image, and interpreting it within the
broader investigative context.

The results show that aggregate accuracy alone is insufficient for determining
whether an automated system provides reliable operational support. Two systems
with similar accuracy may distribute their errors differently between missed
target images and additional false-positive review workload. Conversely,
unchanged or increased aggregate accuracy may coexist with newly introduced
sample-level errors.

The direction of an error is therefore a primary forensic consideration. A
false negative on the \texttt{weapon} class may prevent potentially relevant
material from receiving appropriate review priority. A false positive may
increase the volume of images requiring manual examination, reduce the
selectivity of the triage process, and dilute analyst attention.

Neither error direction should be evaluated in isolation from the intended
workflow. An investigative context in which missed firearm-oriented content is
considered particularly consequential may favor a more permissive operating
configuration. A workflow characterized by extremely large evidence volumes
may instead require greater selectivity to maintain a manageable review burden.
The appropriate balance is therefore case- and task-dependent.

The same distinction applies to the separate \gls{ood} branch. A
\texttt{weapon} assignment on an \gls{ood} image is not a supervised binary
false positive because the input has no
\texttt{weapon}/\texttt{non\_weapon} reference assignment. It may nevertheless
affect the workflow by triggering a flag, changing the priority assigned to an
item, or increasing the amount of material presented for examination.

The \gls{ood} results demonstrate that closed-set systems may produce
operationally consequential assignments even when an input lies outside the
nominal binary task. Explicit management of anomalous, ambiguous, degraded,
synthetic, or semantically adjacent content is therefore necessary. Such
management may involve separate reporting, analyst-visible uncertainty
indicators, rejection mechanisms, or explicit unknown states where supported
by the system.

Where maximum predicted-class probabilities are available, they must be
interpreted only within the corresponding proxy model. They are not calibrated
across architectures and do not independently establish correctness or
forensic reliability. The selected Sigma-Zero case demonstrates that an
incorrect output may be associated with a maximum predicted-class probability
of 1.000.

The commercial software outputs present a further interpretive limitation.
Tags, bookmarks, exported percentages, classification labels, and semantic
retrieval membership do not constitute homogeneous uncertainty measures. They
must not be converted into a common cross-system probability scale unless such
a conversion is independently justified and validated.

Configuration must be treated as part of the evaluated forensic system rather
than as an incidental technical setting. The
\gls{excirefoto2025} sensitivity analysis shows that changing the
\texttt{Distance limit} materially alters recall, false-positive workload, and
\gls{ood} weapon assignment behavior.

The same principle applies to the positive-output mapping adopted for
\gls{griffeye}/\gls{t3kcore}, the exported classification labels used for
\gls{cellebriteinseyets}, and the \texttt{Possible weapons} tag evaluated for
\gls{magnetaxiomtool}/\gls{magnetai}. A reported result is inseparable from the
software version, enabled processing functions, category or query
configuration, output format, and normalization procedure through which it was
obtained.

Operational use therefore requires explicit configuration management. Each
automated output should remain traceably associated with the source item and
its integrity record, together with the software or model version, the active
checkpoint or operational configuration, and the processing functions enabled
at the time of analysis. The record should also preserve the category,
bookmark, label, or semantic query used to obtain the result, the observable
output produced by the system, and the normalization rule applied to that
output. Finally, the experimental or operational condition under which the item
was processed should remain documented together with the subsequent analyst
review and final interpretation.

Software and model updates should also be treated as changes to the evaluated
system. A new release may modify preprocessing, model weights, thresholds,
categories, retrieval behavior, export structures, or default settings.
Results obtained from one frozen version must therefore not be assumed to
remain valid after an update without additional verification.

Traceability is consequently a substantive forensic requirement rather than an
administrative addition. The \gls{fairlab} implementation preserves the links
among source images, derived artifacts, bundle identifiers, proxy-model
predictions, commercial software outputs, normalized decisions, metrics, and
reporting assets.

This structure supports the reconstruction of the analytical process and helps
distinguish system behavior from data-processing errors, incomplete exports,
incorrect mappings, or missing records. An automated result whose provenance
cannot be reconstructed or independently reviewed has limited analytical value
within a forensic workflow.

The residual embedded-metadata analysis reinforces this requirement. Opaque
file names and directory structures reduce direct disclosure of experimental
assignments, but they do not guarantee byte-level blindness. Embedded metadata
may remain available to software systems even when it is not visible in the
normalized result.

The leave-out analysis shows that the identified sensitive-term cases do not
materially alter the aggregate conclusions. It does not establish that the
metadata was never accessed or used. Blind-input claims should therefore
specify whether blindness applies to file-system organization, embedded
metadata, image pixels, or all of these channels.

The qualitative \gls{xai} analysis provides an additional form of diagnostic
transparency for \gls{efficientnetb0}. It preserves a link among the input,
prediction, maximum predicted-class probability, attribution target, and
visualization. This information may support analyst inspection of selected
cases.

The attribution maps do not identify causal features, certify prediction
correctness, or explain the selected black-box systems. They also do not remove
the need to evaluate the model through quantitative error metrics. Explainability
must remain supplementary to robustness testing and expert assessment.

These considerations support a \textit{human-in-the-loop} model of use in which
the analyst remains responsible for interpreting automated outputs within the
investigative context and for reviewing both flagged content and, where
appropriate, material that has not been prioritized by the automated process.
An unflagged item may still retain investigative relevance, and automated
decisions should therefore not be treated as substitutes for expert review.

The analyst should also document the software or model version, the active
configuration, and the normalization rule associated with each result, while
assessing the operational consequences of false negatives, false positives, and
\gls{ood} assignments. Automated outputs should be corroborated with the
remaining forensic evidence, and the final analytical and evidentiary
significance of the material must remain the responsibility of the qualified
human examiner.

The experimental findings therefore do not argue against the use of
\gls{ai}-assisted image analysis in \gls{df}. They define the conditions under
which such use is more defensible: task-specific validation, explicit threat
modeling, traceable configurations, version control, human oversight, and
continued visibility of the system's limitations.


% ─────────────────────────────────────────────────────────────────────────────
\section{Limitations}
\label{sec:conclusions-limitations}
% ─────────────────────────────────────────────────────────────────────────────

The findings apply within the scope of the frozen protocol and the defined case
study. They must not be interpreted as universal validation of the evaluated
models or software systems. The principal limitations are summarized below.

\begin{itemize}

\item \textbf{Dataset size and intended scope.}

The frozen dataset contains 1\,500 manually selected source images: 500
\texttt{weapon} images, 500 \texttt{non\_weapon} images, and 500 images
reserved for the separate \gls{ood} branch. This population supports a
controlled operational robustness study but does not represent the complete
diversity of content encountered in real forensic collections.

\item \textbf{Artificial binary balance.}

The binary subset contains equal numbers of \texttt{weapon} and
\texttt{non\_weapon} images. This balance facilitates controlled comparison of
the two error directions but does not estimate the prevalence of firearm-related
content in operational evidence. Accuracy, precision, and review workload may
change substantially under different class distributions.

\item \textbf{Single-reviewer semantic validation.}

The final operational assignments derive from a documented
\textit{human-in-the-loop} selection procedure conducted by one reviewer. No
independent second annotation or inter-rater agreement measurement was included.
The assignments are traceable but remain dependent on the adopted semantic
criteria and reviewer judgment.

\item \textbf{Source-distribution imbalance.}

The source population is heterogeneous but not uniformly distributed across
the final operational classes. In particular, the
\texttt{non\_weapon} branch is predominantly derived from the
\texttt{05\_deepweb} source group. Source-specific visual regularities may
therefore influence some results despite source-aware fold construction.

\item \textbf{Heterogeneous \gls{ood} population.}

The \gls{ood} branch intentionally combines borderline, ambiguous, degraded,
synthetic, non-firearm weapon imagery, and semantically external content. It is
not a homogeneous supervised class or a universal representation of
out-of-distribution imagery. The reported rates characterize behavior only on
the frozen mixture used in this study.

\item \textbf{Derived-sample dependence.}

The 5\,000 adversarial and 5\,000 anti-forensic artifacts are derived from the
same 1\,000 clean binary images. The 11\,000-item binary result population
therefore represents source--condition combinations rather than 11\,000
independent underlying cases.

Bundle-weighted metrics must not be interpreted as estimates derived from an
independent operational population.

\item \textbf{Finite perturbation suite.}

The adversarial and anti-forensic branches include a finite set of conditions
with one frozen parameterization for each attack or transformation. Different
budgets, severities, optimization procedures, encodings, implementations, or
sequential combinations may produce different robustness profiles.

\item \textbf{Restricted direct-target adversarial evaluation.}

The four model-dependent adversarial attacks were generated against
fold-specific \gls{efficientnetb0} checkpoints. Their evaluation on
\gls{efficientnetb0} constitutes direct-target assessment.

The results obtained on \gls{resnet18}, the \gls{clip}-based proxy, and the
selected black-box software configurations measure only empirical transfer of
the \gls{efficientnetb0}-oriented artifacts. They do not establish robustness
against attacks optimized independently for those systems.

\item \textbf{Artifact-representation effects in black-box transfer.}

The \gls{fgsm}, \gls{sigmazero}, and \gls{superdeepfool} artifacts were
exported as \(224 \times 224\) \texttt{PNG} files, whereas the clean source
population retained heterogeneous original dimensions and formats. The
\gls{onepixel} outputs also used \texttt{PNG} serialization while preserving the
original spatial dimensions.

Consequently, black-box responses to these artifacts may include effects of
resizing or serialization in addition to the optimized perturbation. The
reported transfer results characterize the complete frozen artifacts submitted
to each system and do not isolate the causal contribution of the adversarial
operation alone.

\item \textbf{Proxy-model role.}

\gls{efficientnetb0}, \gls{resnet18}, and the \gls{clip}-based model are
transparent experimental references. They are not reconstructions or validated
surrogates of the internal models used by the selected proprietary software
systems.

Their gradients, probability distributions, and attribution results must not be
used to infer the internal behavior of the black-box systems.

\item \textbf{Probability and score interpretation.}

The maximum predicted-class probabilities produced by the proxy models are
architecture-specific and are not calibrated across models. Exported black-box
percentages, tags, labels, bookmarks, and retrieval results have different
semantics and cannot be treated as directly comparable probabilities.

\item \textbf{Different \gls{ood} prediction units.}

Each proxy-model \gls{ood} rate aggregates 2\,500 fold-specific responses,
corresponding to five checkpoints applied to each of 500 unique images. Each
black-box configuration produces one normalized decision per image.

The proportions describe parallel operating tendencies but are not based on
identical numbers or types of prediction events.

\item \textbf{Black-box version and configuration dependence.}

The black-box results are specific to the evaluated software versions, enabled
functions, imported population, export procedures, and normalization rules.
Internal updates or configuration changes may modify the observed behavior.

The results must not be extrapolated automatically to future releases,
different categories, alternative case exports, or other media-analysis
settings.

\item \textbf{Operational normalization of heterogeneous outputs.}

The common \texttt{weapon}/\texttt{non\_weapon} representation is an explicit
operational recoding of observable outputs. It does not establish that the
selected systems implement equivalent internal binary classification tasks.

The results depend on the \texttt{Possible weapons} tag adopted for
\gls{magnetai} and on the conservative
\texttt{CORE/Violence/Firearm} bookmark mapping used for
\gls{griffeye}/\gls{t3kcore}.

They also depend on the eight firearm-oriented semantic queries and the
\texttt{Distance limit} adopted for \gls{excirefoto2025}, as well as on the
exported \texttt{Armi}, \texttt{Pistola}, and \texttt{Fucile} labels used for
\gls{cellebriteinseyets}.

Alternative defensible mappings would define different operating tasks and may
produce different error profiles.

\item \textbf{Qualitative \gls{xai} scope.}

The retained \gls{xai} analysis is limited to five purposively selected
\gls{efficientnetb0} cases. The set is not statistically representative and
cannot be used to estimate the prevalence of the illustrated behaviors.

The historical \gls{ig} assets were generated using 32 integration steps and a
zero baseline in the preprocessed input space. Numerical convergence was not
established for the retained figures. They are therefore used only as
qualitative diagnostic visualizations and cannot be extended to the black-box
software systems.

\item \textbf{Embedded-metadata sensitivity.}

The blind bundle concealed experimental information through opaque file names,
paths, and directory structure, but it preserved metadata embedded within the
image files. Fifteen files contain predefined sensitive-term matches.

The leave-out analysis identifies limited aggregate variation but does not
establish whether a black-box system accessed or used the embedded fields. A
paired metadata-neutralized bundle was not included.

\item \textbf{Controlled-data and public-artifact boundary.}

Complete source-image collections, generated image trees, proprietary export
contents, and commercial prediction-level outputs cannot all be redistributed
through the public repository because of legal, ethical, institutional,
licensing, and proprietary restrictions.

The repository preserves scripts, manifests, configuration records,
cryptographic digests, normalized summaries, metrics, selected reporting
assets, and validation reports. This supports procedural reproducibility and
artifact-level auditability, but full public re-execution requires authorized
access to the controlled datasets, model checkpoints, and proprietary outputs.

\item \textbf{Legal and evidentiary scope.}

The study evaluates technical and operational triage behavior. It does not
determine the legal admissibility, probative value, or evidentiary weight of an
automated classification.

Such an assessment would require consideration of the applicable legal
framework, chain-of-custody requirements, independent validation, disclosure
obligations, and case-specific evidentiary context.

\end{itemize}

These limitations delimit the claims supported by the experiment. The thesis
provides a protocol-specific assessment of operational robustness. It does not
certify any software product, establish a universal ranking of the evaluated
systems, or constitute a complete adversarial-security assessment.

The study also does not validate a production-grade firearm detector, provide
evidentiary validation of automated classifications, or establish a basis for
autonomous operational deployment.

The findings must therefore be interpreted as evidence about the observable
behavior of the frozen systems, configurations, mappings, dataset, and threat
model evaluated in this thesis.


% ─────────────────────────────────────────────────────────────────────────────
\section{Future Work}
\label{sec:conclusions-future-work}
% ─────────────────────────────────────────────────────────────────────────────

Future research directions derive directly from the experimental findings and
from the limitations of the frozen protocol. The main opportunities for
extension concern the experimental population, adversarial and anti-forensic
threat models, \gls{ood} handling, black-box configuration management,
explainability, operational validation, and independent replication.

A first direction concerns the size and diversity of the experimental
population. Future studies should include larger multi-source collections,
additional firearm-related subcategories, more varied realistic negative
samples, and visual material obtained from different devices, acquisition
procedures, messaging platforms, social networks, and media-processing
pipelines. Increasing source diversity would provide a stronger basis for
assessing whether the operating profiles observed in this study persist across
different acquisition and content-generation conditions.

The class distribution should also be varied. The balanced binary population
used in this thesis, comprising 500 \texttt{weapon} and 500
\texttt{non\_weapon} images, supports controlled comparison of false negatives
and false positives but does not represent the prevalence expected in
operational evidence collections. Evaluations conducted under different class
ratios would make it possible to examine how prevalence affects precision,
false-positive review workload, missed-target risk, and the practical
usefulness of automated triage.

The semantic-validation procedure could be strengthened through the
involvement of multiple independent reviewers. A multi-annotator protocol,
supported by formal agreement measures and an explicit adjudication process,
would quantify uncertainty in borderline cases and reduce dependence on the
judgment of a single reviewer. This would be particularly relevant for
ambiguous firearm-related material and for the heterogeneous
\gls{ood} population.

The \gls{ood} branch itself should be decomposed into more specific
subpopulations. Borderline firearm-related content, non-firearm weapons,
synthetic imagery, degraded media, video-game content, large-scale military
equipment, and semantically unrelated material may produce substantially
different closed-set behaviors. Analyzing these groups separately would help
identify whether the observed weapon-assignment rates are concentrated in
particular semantic or visual categories.

Future work should also examine mechanisms designed explicitly for inputs
outside the nominal binary task. Possible approaches include calibrated
model-specific probability estimates, abstention thresholds, dedicated
\gls{ood} detectors, open-set recognition, and selective-classification
procedures capable of exposing an unknown or ambiguous state to the analyst.

The value of these mechanisms should not be assessed exclusively through
conventional classification metrics. Their effects on missed detections,
false-positive workload, deferred cases, analyst review time, and final
decision quality should also be measured. An effective rejection mechanism,
for example, may reduce inappropriate closed-set assignments while increasing
the number of cases requiring direct human review. This tradeoff should be
treated as an operational property rather than solely as a model-performance
measure.

The adversarial branch should be extended through matched direct-target attacks
against each transparent proxy architecture. In the present protocol,
model-dependent attacks were generated only against fold-specific
\gls{efficientnetb0} checkpoints. Comparable attacks generated independently
against \gls{resnet18} and the \gls{clip}-based proxy would permit a more
systematic distinction between direct-target vulnerability and
cross-architecture transferability.

Such an extension should use harmonized access assumptions and perturbation
budgets where technically meaningful while preserving the distinction among
white-box, score-based, transfer, and model-agnostic conditions. This would
allow attack-specific effects to be separated more clearly from
architecture-specific robustness.

Black-box transfer studies should additionally include clean controls matched
to the adversarial artifacts in spatial resolution and serialization format.
Comparing each manipulated population with its representation-matched clean
control would separate perturbation effects from resizing and encoding effects.

Future adversarial evaluations could also include broader parameter sweeps,
alternative optimization objectives, targeted attacks, adversarial patches,
physical or hybrid manipulations, and conditions designed to remain effective
after ordinary media handling. These experiments would help determine whether
the observed vulnerabilities persist under more varied attack objectives and
delivery mechanisms.

The evaluation of defensive strategies represents a complementary direction.
Potential approaches include adversarial training, robust preprocessing, input
randomization, ensemble decision procedures, and anomaly-detection mechanisms.
Their assessment should consider not only recovery of aggregate accuracy but
also changes in false negatives, false positives, \gls{ood} behavior,
computational cost, and analyst workload. A defense that improves aggregate
performance while materially increasing missed-target risk would not
necessarily constitute an operational improvement.

The anti-forensic branch should likewise be expanded through multiple severity
levels and sequential transformation chains. The current protocol evaluates
one frozen parameterization for each transformation. Future studies should
consider repeated JPEG encoding, platform-specific recompression, multiple
resize ratios, cropping, partial occlusion, screenshot acquisition, and
color-space conversion.

Additional conditions could include stronger and weaker levels of blurring,
combined tonal and geometric transformations, and sequential adversarial and
anti-forensic manipulations. Such experiments would make it possible to
determine whether individually limited effects accumulate when images pass
through realistic evidence-handling, content-sharing, or concealment
processes.

The black-box evaluation should be repeated across additional software
versions, processing modules, operational categories, and documented
configurations. Proprietary systems may change their internal models,
preprocessing procedures, output taxonomies, thresholds, retrieval mechanisms,
or export formats between releases. Findings obtained from one frozen version
should therefore not be assumed to remain stable after software updates.

Longitudinal evaluation would help distinguish behavior that persists across
versions from findings that depend on a particular release or configuration.
It could also support the development of explicit change-management procedures
for operational environments in which software updates may invalidate previous
validation results.

Alternative normalization mappings should be evaluated whenever the observable
outputs support more than one defensible operational interpretation. Such
experiments could consider broader bookmark combinations for
\gls{griffeye}/\gls{t3kcore}, alternative unions of exported categories for
\gls{cellebriteinseyets}, and additional observable tags exposed by
\gls{magnetai}.

For \gls{excirefoto2025}, different firearm-oriented semantic-query sets and
\texttt{Distance limit} values should also be examined systematically. The
objective would not be to identify a universally optimal mapping, but to
quantify how configuration and normalization choices alter recall,
selectivity, \gls{ood} assignment behavior, and human-review workload.

A dedicated controlled experiment should address the residual embedded-metadata
channel identified in the present study. Matched metadata-preserving and
metadata-neutralized copies of the same images should be created while keeping
their visual content unchanged. Both populations could then be processed under
identical software configurations.

A paired comparison would permit direct measurement of output changes
associated with metadata removal and would provide stronger evidence regarding
whether a system accesses or relies on embedded non-pixel information. This
would move beyond the leave-out sensitivity analysis used in the current
protocol and allow a causal effect of metadata modification to be investigated
more directly.

The explainability component should also be strengthened. The five historical
\gls{ig} cases should be regenerated using the hardened adaptive
convergence-validation procedure described in
Section~\ref{subsec:implementation-explainability}. Future analyses should
adopt explicit numerical convergence criteria, multiple attribution baselines,
and larger, systematically sampled case populations.

Different attribution methods could also be compared, together with the
stability of explanations across checkpoints and across clean and manipulated
versions of the same source image. Quantitative measures of attribution
consistency may provide an additional basis for evaluating the stability of
the resulting visualizations.

Explainability should nevertheless remain supplementary to quantitative
robustness testing. Agreement among attribution methods or visual stability
across conditions must not be interpreted as proof of classification
correctness, causal feature relevance, or forensic validity.

A further extension should evaluate the practical consequences of automated
outputs through analyst-centered studies. The present thesis infers
operational relevance from error direction, assignment behavior, and expected
review workload. Controlled user studies could measure these effects directly
by comparing forensic workflows with and without automated triage.

Such studies could measure analyst review time, the number of images examined,
the number of relevant items detected or missed, and the frequency with which
analysts override automated outputs. They could also examine whether false
positives affect attention allocation and whether false negatives influence
final case-level findings.

The practical usefulness of uncertainty indicators, explicit \gls{ood} states,
and explainability information could likewise be evaluated. This would connect
technical robustness metrics more directly to the actual quality of support
provided to the forensic analyst.

The \gls{fairlab} protocol could also be applied to other forensic visual
categories. Potential domains include illicit substances, documents,
screenshots, explicit content, extremist material, vehicles, faces, logos, and
other categories of investigative interest. Each extension would require its
own task definition, semantic-validation criteria, threat model, output
normalization, and operational interpretation.

Finally, the methodological approach would benefit from independent
replication and broader artifact availability. Where licensing, privacy, legal,
or institutional constraints prevent complete redistribution, future
repositories should continue to preserve source and artifact schemas,
executable scripts, configuration records, software-version information,
manifests, and stable identifiers.

Cryptographic digests, sanitized normalized outputs, aggregate and
condition-level metrics, and the corresponding validation and compilation
reports should likewise remain available whenever possible. These artifacts
would support independent verification even when complete datasets or
proprietary exports cannot be distributed publicly.

Independent application of the protocol across laboratories, software
versions, datasets, and forensic tasks would provide stronger evidence
regarding the generalizability of the methodological approach. Such
replication should continue to distinguish procedural reproducibility from
unrestricted access to controlled, licensed, or proprietary material.


% ─────────────────────────────────────────────────────────────────────────────
\section{Closing Remarks}
\label{sec:conclusions-closing}
% ─────────────────────────────────────────────────────────────────────────────

This thesis has shown that the operational robustness of \gls{ai}-based image
classification and semantic retrieval in \gls{df} cannot be characterized
through clean accuracy alone. Systems that perform well under nominal
conditions may still exhibit missed detections, increased false-positive
review workload, non-zero \gls{ood} weapon-assignment rates, configuration
sensitivity, or substantial degradation under specific manipulated
conditions.

A central contribution of the study is the parallel evaluation of transparent
proxy models and selected black-box software configurations on the same frozen
source population and derived perturbation suite. This design combines
experimental control with operational relevance while maintaining a clear
distinction between inspectable model behavior and the normalized observable
outputs of proprietary software systems.

The \gls{fairlab} protocol provides the traceable methodological structure
through which this evaluation is conducted. It connects dataset construction,
human semantic validation, deterministic partitioning, proxy-model evaluation,
perturbation generation, blind-bundle construction, black-box output
normalization, quantitative analysis, qualitative explainability, and
integrity controls.

Within the defined experimental scope, no evaluated system or configuration
minimizes all dimensions of operational risk. High target recall may coexist
with substantial false-positive workload, while strong clean accuracy may
coexist with severe direct-target adversarial vulnerability. More selective
configurations may reduce review workload while increasing missed detections,
whereas more permissive configurations may reduce false negatives at the cost
of routing larger quantities of unrelated material to human review.

The separate \gls{ood} analysis further shows that closed-set systems may
assign out-of-task content to an operational category even when no supervised
binary reference assignment exists. Such behavior reinforces the need to
distinguish \gls{ood} assignment from ordinary false-positive error and to make
anomalous or ambiguous content visible to the analyst.

These findings do not preclude the use of \gls{ai}-assisted software in
forensic workflows. They indicate that such use is more defensible when
supported by task-specific validation on representative data, an explicit
threat model, documented software and model versions, controlled
configurations, and reproducible output mappings.

Source-to-output traceability must also be preserved, while false negatives,
false positives, \gls{ood} assignments, and configuration-dependent behavior
must remain visible rather than being concealed behind a single aggregate
metric. Automated classifications, tags, bookmarks, and semantic retrieval
results should therefore be interpreted as triage indicators rather than as
autonomous evidentiary conclusions.

Such outputs may assist in organizing and prioritizing large evidence
collections, but they do not replace contextual analysis, corroboration,
chain-of-custody requirements, or expert judgment. Maximum predicted-class
probabilities and other system-specific scores likewise cannot independently
establish correctness or forensic reliability.

Operational robustness must therefore be demonstrated under the conditions in
which a system is expected to operate. It cannot be presumed from nominal
performance, proprietary status, or the apparent confidence of an automated
output.

In forensic image analysis, automation remains most defensible when it operates
as a controlled, auditable, traceable, and reviewable component of a
\textit{human-in-the-loop} workflow.
